\documentclass[12pt,a4paper]{article}

\usepackage{fontspec}
\setmainfont{Times New Roman}
\setsansfont{PT Sans}
\setmonofont{PT Mono}

\usepackage{polyglossia}
\setmainlanguage{russian}
\setotherlanguage{english}

\usepackage{geometry}
\geometry{left=30mm,right=15mm,top=20mm,bottom=20mm}

\usepackage{setspace}
\onehalfspacing

\usepackage{tikz}
\usetikzlibrary{arrows.meta,positioning,calc}
\usepackage{float}
\usepackage{caption}
\usepackage{hyperref}
\usepackage{fancyhdr}

\hypersetup{
    colorlinks=true,
    urlcolor=blue,
    linkcolor=black,
    pdftitle={Информационные системы. Лабораторная работа 1},
    pdfauthor={Кузьмин Дмитрий Анатольевич}
}

\pagestyle{fancy}
\fancyhf{}
\fancyfoot[C]{\thepage}
\renewcommand{\headrulewidth}{0pt}
\setlength{\parskip}{0.4em}
\setlength{\parindent}{1.25cm}
\setlength{\emergencystretch}{3em}

\captionsetup{font=small,skip=6pt,labelsep=endash}

\tikzset{
    uml/.style={
        draw,
        inner sep=3pt,
        font=\scriptsize\linespread{1.05}\selectfont,
        fill=white,
        align=center
    },
    dep/.style={-{Stealth[length=2.2mm]}, semithick},
    assoc/.style={semithick},
    use/.style={-{Stealth[length=2.2mm,open]}, dashed, semithick}
}

\newcommand{\cls}[3]{%
    \begin{tabular}{@{\hspace{3pt}}l@{\hspace{3pt}}}
        \multicolumn{1}{c}{\textbf{#1}}\\
        \hline
        #2\\
        \hline
        #3
    \end{tabular}%
}

\newcommand{\pkg}[2]{%
    \begin{tabular}{@{\hspace{3pt}}l@{\hspace{3pt}}}
        \multicolumn{1}{c}{\textbf{#1}}\\
        \hline
        #2
    \end{tabular}%
}

\newcommand{\enm}[2]{%
    \begin{tabular}{@{\hspace{3pt}}l@{\hspace{3pt}}}
        \multicolumn{1}{c}{\textbf{<<enumeration>>}}\\
        \multicolumn{1}{c}{\textbf{#1}}\\
        \hline
        #2
    \end{tabular}%
}

\newcommand{\clsa}[2]{%
    \begin{tabular}{@{\hspace{3pt}}l@{\hspace{3pt}}}
        \multicolumn{1}{c}{\textbf{#1}}\\
        \hline
        #2
    \end{tabular}%
}

\begin{document}

\begin{titlepage}

\begin{center}
    \LARGE Национальный исследовательский университет ИТМО\\[8mm]

    Факультет Программной Инженерии и Компьютерной техники\\[37mm]

    \huge \textbf{Информационные системы}\\[2mm]
    \textbf{Лабораторная работа 1} \\[2mm]



\end{center}

\vfill




\begin{flushright}

    \large Выполнил:

    Кузьмин Дмитрий Анатольевич\\[3mm]

    Группа:

    P3309\\[3mm]

    Преподаватель:

    Николаев Владимир Вячеславович
\end{flushright}


\vfill

\begin{center}
    Санкт-Петербург 2026
\end{center}

\end{titlepage}

\section{UML-диаграммы классов и пакетов}

Предметная область работы~--- учёт драконов. Серверная часть написана на Java~17 с использованием Jakarta~EE~10 (JAX-RS, CDI, Bean Validation, JPA/Hibernate) и поставляется архивом \texttt{lab1.war}. Данные хранятся в PostgreSQL. Клиент~--- статические страницы, которые вызывают REST API по пути \texttt{/api} и получают сигнал об изменениях через WebSocket \texttt{/ws/dragons}.

\subsection{Диаграмма пакетов}

Классы разложены по пакетам в соответствии со слоями. Пакет \texttt{lab.rest} принимает HTTP-запросы. Пакет \texttt{lab.service} содержит сценарии и проверки. Пакет \texttt{lab.repository} обращается к базе. Пакет \texttt{lab.entity} описывает сохраняемые объекты. Пакет \texttt{lab.dto} задаёт тела запросов и ответов и не зависит от сущностей. Пакет \texttt{lab.mapper} собирает \texttt{DragonResponse} из сущности. Пакет \texttt{lab.ws} рассылает клиентам уведомление после изменений.

Стрелка на рисунке~\ref{fig:packages} означает зависимость: пакет на стороне начала стрелки импортирует классы пакета, на который она указывает.

\begin{figure}[H]
\centering
\begin{tikzpicture}[node distance=8mm]
    \node[uml] (rest) {\pkg{lab.rest}{
        RestApplication\\
        JsonConfig\\
        DragonResource\\
        SpecialResource\\
        IllegalArgumentMapper\\
        MissingExceptionMapper
    }};
    \node[uml, below left=14mm and -2mm of rest] (service) {\pkg{lab.service}{
        DragonService\\
        SpecialService\\
        MissingException
    }};
    \node[uml, below=of service] (repo) {\pkg{lab.repository}{
        DragonRepository\\
        DragonPage\\
        SpecialRepository\\
        DbSetup\\
        DbSetupStarter
    }};
    \node[uml, below=of repo] (entity) {\pkg{lab.entity}{
        Dragon, Coordinates, DragonCave\\
        DragonHead, Person, Location\\
        KillerTeam\\
        Color, Country, DragonCharacter
    }};
    \node[uml, right=16mm of service] (dto) {\pkg{lab.dto}{
        DragonRequest, DragonResponse\\
        PageResponse, OptionsResponse\\
        ErrorResponse\\
        TeamRequest, TeamResponse\\
        SendTeamRequest\\
        CoordinatesOption, CaveOption\\
        HeadOption, PersonOption\\
        LocationOption
    }};
    \node[uml, right=16mm of repo] (mapper) {\pkg{lab.mapper}{
        DragonMapper
    }};
    \node[uml, right=16mm of rest] (ws) {\pkg{lab.ws}{
        DragonSocket
    }};

    \draw[dep] (rest) -- (service);
    \draw[dep] (rest.east) -- (ws.west);
    \draw[dep] (rest.south east) -- (dto.north west);
    \draw[dep] (service) -- (repo);
    \draw[dep] (service.east) -- (dto.west);
    \draw[dep] (service.south east) -- (mapper.north west);
    \draw[dep] (repo) -- (entity);
    \draw[dep] (mapper.south) -- (entity.east);
    \draw[dep] (mapper.north) -- (dto.south);
\end{tikzpicture}
\caption{Диаграмма пакетов серверной части}
\label{fig:packages}
\end{figure}

\subsection{Диаграмма классов предметной области}

На рисунке~\ref{fig:domain} показаны сущности JPA и перечисления. Сплошная линия~--- ассоциация, подписи задают кратность. Пунктир~--- использование перечисления. Общие объекты (\texttt{Coordinates}, \texttt{DragonCave}, \texttt{DragonHead}, \texttt{Person}, \texttt{Location}) могут быть связаны с несколькими драконами или людьми, поэтому связь не является композицией.

\begin{figure}[H]
\centering
\resizebox{\linewidth}{!}{%
\begin{tikzpicture}[mul/.style={font=\tiny, fill=white, inner sep=1.2pt}]
    \node[uml] (coord) at (-5.6,0) {\clsa{Coordinates}{
        id: Integer\\
        x: float\\
        y: Integer
    }};
    \node[uml] (cave) at (0,0) {\clsa{DragonCave}{
        id: Integer\\
        depth: double\\
        numberOfTreasures: Long
    }};
    \node[uml] (head) at (8.6,0) {\clsa{DragonHead}{
        id: Integer\\
        size: long
    }};

    \node[uml] (character) at (-5.6,-4.8) {\enm{DragonCharacter}{
        WISE\\
        GOOD\\
        CHAOTIC\\
        FICKLE
    }};
    \node[uml] (dragon) at (0,-4.8) {\cls{Dragon}{
        id: Integer\\
        name: String\\
        creationDate: LocalDate\\
        age: long\\
        wingspan: float\\
        weight: Float\\
        character: DragonCharacter
    }{fillCreationDate()}};
    \node[uml] (team) at (5.4,-4.8) {\clsa{KillerTeam}{
        id: Integer\\
        name: String\\
        cave: DragonCave\\
        members: Set<Person>
    }};

    \node[uml] (location) at (-5.6,-10.2) {\clsa{Location}{
        id: Integer\\
        x: double\\
        y: double\\
        z: float\\
        name: String
    }};
    \node[uml] (person) at (0,-10.2) {\clsa{Person}{
        id: Integer\\
        name: String\\
        eyeColor: Color\\
        hairColor: Color\\
        birthday: LocalDateTime\\
        nationality: Country\\
        location: Location
    }};
    \node[uml] (color) at (5.6,-9.2) {\enm{Color}{
        GREEN\\
        BLACK\\
        BLUE\\
        YELLOW\\
        ORANGE
    }};
    \node[uml] (country) at (5.6,-12.6) {\enm{Country}{
        RUSSIA\\
        USA\\
        GERMANY\\
        SPAIN\\
        INDIA
    }};

    \draw[assoc] (dragon.north) -- (cave.south)
        node[mul, pos=0.15] {*} node[mul, pos=0.85] {1};
    \draw[assoc] (dragon.north west) -- (coord.south)
        node[mul, pos=0.18] {*} node[mul, pos=0.82] {1};
    \draw[assoc] (dragon.north east) -- (head.south)
        node[mul, pos=0.18] {*} node[mul, pos=0.82] {0..1};
    \draw[use] (dragon.west) -- (character.east);

    \draw[assoc] (dragon.south) -- (person.north)
        node[mul, pos=0.15] {*} node[mul, pos=0.85] {0..1};
    \draw[assoc] (person.west) -- (location.east)
        node[mul, pos=0.15] {*} node[mul, pos=0.85] {1};
    \draw[use] (person.east) -- (color.west);
    \draw[use] (person.south east) -- (country.west);

    \draw[assoc] (team.north west) -- (cave.east)
        node[mul, pos=0.22] {*} node[mul, pos=0.78] {0..1};
    \draw[assoc] (team.south west) -- (person.north east)
        node[mul, pos=0.2] {*} node[mul, pos=0.8] {*};
\end{tikzpicture}%
}
\caption{Диаграмма классов предметной области}
\label{fig:domain}
\end{figure}

\texttt{Dragon} обязательно ссылается на координаты и пещеру. Голова и убийца могут отсутствовать. Дата создания проставляется перед первой записью, если клиент её не передал. \texttt{Person} обязательно имеет местоположение; цвет глаз обязателен, цвет волос и гражданство~--- нет. \texttt{KillerTeam} объединяет убийц связью <<многие ко многим>> и может быть направлена в пещеру.

\subsection{Диаграмма классов приложения}

На рисунке~\ref{fig:app} показаны классы, через которые проходят запросы. Ресурсы делегируют работу сервисам и после создания, изменения, удаления дракона, а также после операций с командой вызывают \texttt{DragonSocket.broadcast()}. Сервисы открывают транзакцию, проверяют данные и обращаются к репозиториям. \texttt{DragonMapper} переводит сущность в ответ API. Если объект не найден, выбрасывается \texttt{MissingException}; обработчики в \texttt{lab.rest} превращают её и ошибки проверки в JSON.

\begin{figure}[H]
\centering
\resizebox{\linewidth}{!}{%
\begin{tikzpicture}
    \node[uml] (mapper) at (6.4,4.6) {\clsa{DragonMapper}{
        $+$ toResponse(dragon): DragonResponse
    }};
    \node[uml] (socket) at (0,-10.6) {\cls{DragonSocket}{
        sessions: Set<Session>
    }{
        $+$ open(session)\\
        $+$ close(session)\\
        $+$ broadcast()
    }};
    \node[uml] (dres) at (0,0) {\cls{DragonResource}{
        $-$ dragons: DragonService
    }{
        $+$ list(...): PageResponse\\
        $+$ options(): OptionsResponse\\
        $+$ get(id): DragonResponse\\
        $+$ create(request): Response\\
        $+$ update(id, request): DragonResponse\\
        $+$ delete(id): Response
    }};
    \node[uml] (dsvc) at (6.4,0) {\cls{DragonService}{
        $-$ dragons: DragonRepository\\
        $-$ validator: Validator
    }{
        $+$ page(...): PageResponse\\
        $+$ get(id): DragonResponse\\
        $+$ create(request): DragonResponse\\
        $+$ update(id, request): DragonResponse\\
        $+$ delete(id)\\
        $+$ options(): OptionsResponse
    }};
    \node[uml] (drepo) at (13.2,0.3) {\cls{DragonRepository}{
        em: EntityManager
    }{
        $+$ search(...): DragonPage\\
        $+$ find(id): Dragon\\
        $+$ persist(entity)\\
        $+$ deleteDragon(dragon)\\
        $+$ allCoordinates(): List\\
        $+$ allCaves(): List\\
        $+$ allPersons(): List\\
        $+$ allHeads(): List\\
        $+$ allLocations(): List
    }};
    \node[uml] (sres) at (0,-5.4) {\cls{SpecialResource}{
        $-$ specials: SpecialService
    }{
        $+$ maxDate(): DragonResponse\\
        $+$ byName(part): List\\
        $+$ byCharacter(character): List\\
        $+$ teams(): List\\
        $+$ createTeam(request): Response\\
        $+$ send(request): TeamResponse
    }};
    \node[uml] (ssvc) at (6.4,-5.4) {\cls{SpecialService}{
        $-$ specials: SpecialRepository\\
        $-$ dragons: DragonRepository
    }{
        $+$ maxCreationDate(): DragonResponse\\
        $+$ byName(part): List\\
        $+$ byCharacter(character): List\\
        $+$ teams(): List\\
        $+$ createTeam(request): TeamResponse\\
        $+$ send(request): TeamResponse
    }};
    \node[uml] (srepo) at (13.2,-5.6) {\cls{SpecialRepository}{
        em: EntityManager
    }{
        $+$ dragonWithMaxCreationDate()\\
        $+$ dragonsByName(part): List\\
        $+$ dragonsByCharacter(c): List\\
        $+$ createTeam(name, ids): int\\
        $+$ sendTeam(teamId, caveId): int\\
        $+$ teams(): List
    }};
    \node[uml] (missing) at (13.2,-10.4) {\clsa{MissingException}{
        message: String
    }};

    \draw[dep] (dres) -- (dsvc);
    \draw[dep] (dsvc) -- (drepo);
    \draw[dep] (dsvc.north) -- (mapper.south);
    \draw[dep] (sres) -- (ssvc);
    \draw[dep] (ssvc) -- (srepo);
    \draw[dep] (ssvc.north east) -- (drepo.south);
    \draw[dep] (ssvc.east) -- ++(0.55,0) |- (mapper.east);
    \draw[dep] (sres.south) -- (socket.north);
    \draw[dep] (dres.west) -- ++(-0.55,0) |- (socket.west);
    \draw[use] (ssvc.south east) -- (missing.north west);
    \draw[use] (dsvc.east) -- ++(0.7,0) |- (missing.north);
\end{tikzpicture}%
}
\caption{Диаграмма классов прикладного слоя}
\label{fig:app}
\end{figure}

Обычный сценарий дракона: страница списка с номером страницы, размером, сортировкой и фильтром по полю; чтение одной записи; создание; изменение; удаление. Справочник \texttt{options} отдаёт уже сохранённые координаты, пещеры, головы, людей и места, чтобы их можно было выбрать повторно. Специальные операции: дракон с наибольшей датой создания, поиск по части имени, поиск по характеру, список команд, создание команды из выбранных убийц и отправка команды в пещеру.

\section{Исходный код}

Исходный код системы расположен в репозитории на ветке \texttt{for\_gelios}:

\begin{center}
\url{https://github.com/Dmitry-Kuzmin456/is_lab_1/tree/for_gelios}
\end{center}

\section{Выводы}

В работе построена информационная система учёта драконов. Модель данных покрывает дракона и связанные с ним координаты, пещеру, голову и убийцу, а также команду убийц, которую можно собрать и направить в пещеру. Ограничения предметной области заданы и в сущностях, и при разборе запроса, поэтому в базу не попадают пустое имя, неположительные возраст, вес, размах крыльев и число сокровищ, неизвестные значения перечислений.

Разделение на пакеты отделяет контракт API от хранения. Ресурсы не собирают SQL и не заполняют сущности, репозитории не знают формат JSON, а объекты \texttt{dto} не тянут за собой таблицы. За счёт этого постраничный список, фильтр и специальные выборки добавляются в своём слое, не переписывая остальные. Общие координаты, пещеры, головы и люди не копируются при каждом драконе: запись ссылается на уже существующий объект или создаёт новый, а при удалении неиспользуемый объект убирается.

Изменяющие операции сообщают открытым страницам о необходимости обновить данные. Так список и экран специальных операций остаются согласованными между несколькими открытыми клиентами без периодического опроса сервера.

\end{document}
